\section{Actual integral extraction of the cubic strict Selmer group}
\label{sec:cubic-selmer-extraction}

Keep the actual cubic seed, ring and complete model of
AC-G49--G51.  Thus $\ell\ne2,3,269$ is in the transported
local range, all old \texttt{GD} conditions and the stated
source inputs remain in force, and the marked element
$g_1$ lies in $H=G_K$ with
$\gamma(g_1)^2-1\in O^\times$.
We additionally retain the full \texttt{LemmaM'} of
\path{CommutativeAlgebra/OrbitVanishingLocal.lean}.
This is the integral cotangent-length input used in the
source's \texttt{model\_lower}; it is not deduced from A.

Write $D=\operatorname{Frac}(O)/O$ and
\[
 \alpha_j=\chi^{-j}\gamma^{-1}|_{G_K}
       =\chi^{-j}\psi^{-1}\kappa_\ell^{-e}|_{G_K}.
                                                               \tag{G52.1}
\]
Define the strict group used in this section explicitly:
\[
 \mathcal S_j=
 \ker\!\left(
 H^1(G_K,D(\alpha_j))\longrightarrow
 \prod_{w\mid\ell}H^1(G_{K_w},D(\alpha_j))
 \times\prod_{w\nmid\ell}H^1(I_w,D(\alpha_j))
 \right).
                                                               \tag{G52.2}
\]
The second product ranges over finite places.
This is the strict coefficient-prime condition and
unramified condition elsewhere.  We do not silently
replace it by a differently defined \texttt{SelAll}.

The exact source convention is
\path{GaloisLibrary/SelmerTransport/ConjStableShaFam.lean},
\texttt{shaFam} and \texttt{ShaU}: the family for a positive
integer $n$ contains all inertia groups and all decomposition
groups above the rational primes dividing $n$.
Consequently (G52.2) is precisely
$\texttt{ShaU}(D(\alpha_j),\ell,\top)$ over $K$.
The source coefficient $\texttt{QO}=O\otimes_{\mathbb Z_\ell}
(\mathbb Q_\ell/\mathbb Z_\ell)$ is canonically $D$:
tensor $0\to\mathbb Z_\ell\to\mathbb Q_\ell
\to\mathbb Q_\ell/\mathbb Z_\ell\to0$ with the finite free
$\mathbb Z_\ell$-module $O$.
For every positive integer $N$, therefore,
\[
 \mathcal S_{j,N}:=
 \texttt{ShaU}(D(\alpha_j),\ell N,\top)
       \ \subseteq\ \mathcal S_j.
\]
The left side imposes additional decomposition-group vanishing
above $N$.  Its inclusion follows by forgetting those extra
conditions, exactly as in \texttt{shaFam\_mono} and
\texttt{V'\_le\_V} of
\path{CharacterLevelBounds/ShaFamMono.lean}.
The same inclusion holds with $N$ replaced by $269N$.
Thus the source's stronger \texttt{ShaU} condition at
$\ell N$ is controlled by the comparison below; it is not
being identified with a group having a relaxed condition at
$\ell$.

\begin{lemma}[AC-G52-normalization: the actual two-dimensional cubic cocycle]
\label{AC-G52-normalization}
Let $L_j=O E_{10}\oplus O E_{21}$ be the weight-$-1$
part of the adjoint lattice of $r_\chi$.
For any discrete $O$-module $M$ on which the induced
action is continuous, identify $M\otimes L_j$ with
column pairs $(x,y)$ and denote that action by $\rho_j$.
For $h\in H$ and the normalized complex conjugation $c$,
\[
 \begin{split}
 \rho_j(h)(x,y)
    &=\bigl(\chi(h)^{-j}\gamma(h)^{-1}x,
             \chi(h)^j\gamma(h)^{-1}y\bigr),\\
 \rho_j(c)(x,y)&=(-y,-x).
 \end{split}                                                   \tag{G52.3}
\]
There is a unique $O$-linear choice of representative
for each class in $H^1(\Gamma,M\otimes L_j)$ satisfying
the source's actual matrix \texttt{NormCoc} conditions.
In particular it is an injective $O$-linear map from
cohomology to normalized continuous cocycles.
It is compatible with every $O$-linear change of $M$.
\end{lemma}
\begin{proof}
The formula on $H$ is obtained by conjugating the two
matrix units by
$\operatorname{diag}(\chi^j\gamma,1,\chi^j\gamma^{-1})$.
The differential of the semidirect involution at $c$ is
$X\mapsto-JX^tJ$, giving the second formula.
It preserves this two-dimensional weight space.

Use the source coboundary sign
$\partial X(g)=X-\rho_j(g)X$.
Let $z$ be a continuous cocycle.
The equation $c^2=1$ gives
$z(c)+\rho_j(c)z(c)=0$, so its two coordinates are equal.
Put $Y=z(c)/2$.  Since $2\in O^\times$,
\[
                     z_1=z-\partial Y,
                     \qquad z_1(c)=0.                       \tag{G52.4}
\]
Every coboundary preserving the last condition is
$\partial X$ with $X=(a,-a)$, because these are exactly
the $c$-fixed vectors.

Set $t=\chi(g_1)$ and $u=\gamma(g_1)$.
Write $z_1(g_1)=(x,y)$ and
$g_1^*=cg_1^{-1}c$.
The cocycle equation, (G52.3) and $z_1(c)=0$ give
\[
                  z_1(g_1^*)=(t^{-j}u\,y,t^ju\,x).
                                                               \tag{G52.5}
\]
Anticyclotomy also gives
$q(g_1^*)=t$ and $\gamma(g_1^*)=u^{-1}$.
The derivative of the source matrix expression
$z(c,g_1)=r(g_1).\mathrm{mat}-r(g_1^*).\mathrm{mat}$
therefore has $(1,0)$ entry
\[
                         \Lambda(z_1)=t^ju\,x-y.             \tag{G52.6}
\]
Its $(0,1)$ entry vanishes identically, because the
cocycle has weight $-1$.
These are precisely the two \texttt{zvar} fields of
\texttt{NormCoc}; the remaining field is (G52.4).

For $X=(a,-a)$, a direct substitution into (G52.3) gives
\[
 \Lambda(\partial X)=\Delta a,
        \qquad\Delta=t^j(u-u^{-1})\in O^\times.              \tag{G52.7}
\]
The unit assertion follows from the old marked condition
$u^2-1\in O^\times$.
Thus put $a=\Delta^{-1}\Lambda(z_1)$ and
\[
                        \nu(z)=z_1-\partial(a,-a).           \tag{G52.8}
\]
This is a continuous normalized representative of the
same cohomology class, and depends $O$-linearly on $z$.

There is no nonzero normalized coboundary.  Indeed its
value at $c$ first forces its defining vector to have
the form $(a,-a)$, and (G52.7) then forces $a=0$.
Consequently two cohomologous normalized cocycles agree.
Formula (G52.8) therefore descends to the asserted
injective section on $H^1$, with all linearity and
coefficient-compatibility assertions.

This replaces the special one-dimensional section
\texttt{shaInj\_norm}/\texttt{shaInj\_mat} in
\path{UnitaryThreeGroups/ShaInjGen.lean} by the actual
cubic two-dimensional section.
It does not assume $\rho_j(g_1)-1$ invertible.
For example $j=1$, $t=u=\omega$ makes its second
diagonal factor $1$, while (G52.7) remains a unit at
$\ell\ne3$.  The marked normalization, rather than
inversion of that matrix, is the valid argument.
\end{proof}

\begin{theorem}[AC-G52: strict cubic length is bounded by the actual model cotangent length]
\label{AC-G52}
Let $\mathcal M_j$ be the actual graded model of AC-G51
for the seed with exponent $j\in\{1,2\}$.
Retaining \texttt{LemmaM'}, one has the exact integral
comparison
\[
                  \operatorname{length}_O\mathcal S_j
                        \le \mathcal M_j.\mathrm{len}.       \tag{G52.9}
\]
Both sides may initially be infinite.  For each
$0\ne b\in O$, the proof gives an actual additive
injection of $\mathcal S_j[b]$ into the model's
weight-$-1$ torsion derivations, with all the integral
lifting properties required by \texttt{LemmaM'}.
There is no extra factor of $2$, Fourier denominator,
or finite-coefficient correction in (G52.9).
In particular the identical upper bound holds for
$\operatorname{length}_O\mathcal S_{j,N}$ for every $N>0$.

This closes the actual integral extraction premise in
the stated coefficient-prime range.  It does not supply
an upper bound for $\mathcal M_j.\mathrm{len}$, prove
\texttt{PtHyp}, or remove the separate Type I and mixed
tangent requirements.
\end{theorem}
\begin{proof}
\proofstep{Identify the actual coefficient and local conditions.}
Replacing the second basis vector in (G52.3) by its
negative turns the action of $c$ into ordinary exchange.
The two $H$-characters are $\alpha_j$ and $\alpha_j^c$.
This gives the explicit integral identification
\[
                  L_j\simeq\operatorname{Ind}_{H}^{\Gamma}
                                      O(\alpha_j).            \tag{G52.10}
\]
Apply the already used integral Shapiro identification
to this rank-two module; no character projector is
inserted.  Restricting it to a rational decomposition
or inertia group decomposes by the corresponding
places of $K$.  The local Shapiro identifications
therefore identify (G52.2) with the subgroup
$\mathcal T_j\subset H^1(\Gamma,D\otimes L_j)$ which
restricts to zero on the coefficient-prime decomposition
group and on every inertia group away from $\ell$.
This use of Shapiro is compatible with the explicit
local conditions, including the ramified places; it
does not replace inertia by decomposition at those
places.

\proofstep{Normalize every torsion class integrally.}
Fix $0\ne b\in O$.
For a class $\zeta\in\mathcal T_j[b]$, apply
AC-G52-normalization with $M=D$.
Because its section $\nu$ is $O$-linear,
\[
                        b\nu(\zeta)=\nu(b\zeta)=0.
                                                               \tag{G52.11}
\]
Thus every value lies in $D[b]\otimes L_j$.
Use the canonical isomorphism
$O/b\simeq D[b]$, $\overline a\mapsto a/b+O$, in
both coordinates to obtain a locally constant matrix
cocycle
\[
 C_\zeta:\Gamma\longrightarrow
       (O/b)E_{10}\oplus(O/b)E_{21}.                          \tag{G52.12}
\]
The map $\zeta\mapsto C_\zeta$ is $O$-linear and
injective.  Its cocycles are the actual normalized
weight-$-1$ cocycles for the cubic basepoint modulo $b$.

Take $g\in\mathcal K_\chi$.
By definition it belongs to inertia away from $\ell$.
The class $\zeta$ is locally zero there, so its
representative is a coboundary on that inertia group.
AC-G49 proves $r_\chi(g)=1$, hence its adjoint action
on $D\otimes L_j$ is the identity.
Its coboundary value at $g$ is therefore zero.
Consequently
\[
                        C_\zeta(g)=0
                           \qquad(g\in\mathcal K_\chi).     \tag{G52.13}
\]
The adjoint action is trivial on these elements, so
the cocycle identity also forces zero on every conjugate
and every product of them.  Continuity then gives the
same killing condition on their closed normal closure,
as needed by the actual ring.

\proofstep{The local cobounding vector is itself $b$-torsion.}
Let $D_\ell$ be the actual chosen rational local
decomposition group.  Its restriction class is zero,
so there is $X\in D\otimes L_j$ with
\[
                  \nu(\zeta)(h)=X-\rho_j(h)X
                                      \qquad(h\in D_\ell).
                                                               \tag{G52.14}
\]
The exact old \texttt{GD} consequence in
\path{UnitaryThreeGroups/DatFacts.lean},
\texttt{dat\_facts}, supplies an element $s$ of local
inertia with $\gamma(s)^2-1\in O^\times$; the actual
local embedding and inertia identification are
\path{UnitaryThreeGroups/DatLoc.lean}, \texttt{dat\_loc}.
Here $\eta(s)=1$ and $\chi(s)=1$, since
$\ell\ne3,269$ and the cubic character is unramified
at the coefficient prime.
Thus $\rho_j(s)=\gamma(s)^{-1}I$ on $L_j$, and
$1-\gamma(s)^{-1}$ is a unit.

Multiply (G52.14) by $b$ and use (G52.11).
At $s$ this gives
$(1-\gamma(s)^{-1})bX=0$, hence $bX=0$.
Under $D[b]\otimes L_j\simeq(O/b)\otimes L_j$,
(G52.14) is therefore a coboundary with a vector in
the finite coefficient module itself.
Together with (G52.12)--(G52.13), this verifies every
field of the cubic version of \texttt{shaCoc} in
\path{UnitaryThreeGroups/ConjStable.lean}.

\proofstep{Construct the actual square-zero global point.}
Put $Q=O/b$ and form the complete local $O$-algebra
$O\oplus Q$, with square-zero second summand.
For $C=C_\zeta$, set
\[
 r_C(g).\mathrm{mat}
      =(I+\varepsilon C(g))r_\chi(g).\mathrm{mat},
             \qquad r_C(g).\mathrm{mult}=\eta(g).             \tag{G52.15}
\]
The cocycle identity says exactly that (G52.15) is a
homomorphism: expand its product and use
$\varepsilon^2=0$ and the differential adjoint action.
It is continuous, has the specified residual
representation and kills the actual inertia kernels
by (G52.13).
The normalized-cocycle equations say exactly that
it is normalized.  This last assertion is the
seed-independent theorem
\path{UnitaryThreeGroups/DefNorm.lean},
\texttt{isDeform\_normalised\_iff}, or follows by
expanding its matrix at $c$ and its two \texttt{zmat}
entries.

Choose a lift $Y\in L_j$ of the finite local vector
from (G52.14).  The representation obtained by
strict conjugation of the integral local basepoint
by $I+\varepsilon Y$, over $O\oplus O$, has cocycle
$Y-\rho_j(h)Y$.
It kills $P_\chi$ by the full strict conjugation
stability proved in AC-G49-inert.
Reduction of its second summand modulo $b$ is exactly
the restriction of (G52.15), by (G52.14).
Thus (G52.15) satisfies the actual local condition,
including over these nonreduced test rings.

The universal property of $S$ now gives its unique
homomorphism $F_C:S\to O\oplus Q$ producing (G52.15).
Its first component is $x_0$ by framed uniqueness.
Write its second component as $D_C:S\to Q$.
This is a derivation at $x_0$, and if
$d_0(g),d_1(g),d_2(g)$ are the diagonal entries of
the cubic basepoint, then
\[
          D_C(r^{\rm univ}(g).\mathrm{mat}_{ij})
                               =d_j(g)C(g)_{ij}.             \tag{G52.16}
\]
Every $d_j(g)$ is an integral unit.
The derivative consequently has weight $-1$ on the
actual entries.  The integral local square-zero
point above supplies its compatible integral local
derivation annihilating $P_\chi$.
Property 5 of AC-G51 therefore produces an integral
weight-$-1$ derivation $\widetilde D_C$ on $A$ lifting
$D_Cq_A$.

This is the entire construction required in the
source's \texttt{main\_der} in
\path{UnitaryThreeGroups/ModelLower.lean}.
The old powers of one character are replaced in
(G52.16) by the actual diagonal units.
Continuity of the integral local coboundary follows
directly from the continuous cubic basepoint and
constant $Y$; no new local constancy hypothesis is
inserted.

\proofstep{An injective additive family of derivations.}
The map $C\mapsto D_C$ is additive.
Indeed a sum of derivations at $x_0$ gives another
map to $O\oplus Q$, and (G52.16) identifies all its
universal matrix entries with those of the point for
the sum of the cocycles.  Framed uniqueness identifies
the two maps.  The same argument proves scalar
compatibility, although only additivity is needed.

The map
\[
 \mathcal S_j[b]\longrightarrow
              \operatorname{Der}_{x}(A,O/b),
       \qquad \zeta\longmapsto D_{C_\zeta}q_A              \tag{G52.17}
\]
is injective.  If its value is zero, take a homogeneous
preimage in $A$ of each actual matrix entry, supplied
by AC-G51.  Formula (G52.16) and the unit $d_j(g)$
force every value of $C_\zeta$ to vanish.
The injectivity in (G52.12) then gives $\zeta=0$.

Each derivation in (G52.17) kills every equation
$F_i$ of weight $-1$, because $q_A(F_i)=0$.
It also admits the integral weight-$-1$ lift just
constructed.  The origin kills $\mathcal A_{-1}$
and all $F_i$ lie in that weight part.
These are all hypotheses of the retained
\texttt{LemmaM'}, applied to $\mathcal S_j[b]$.
Its conclusion is
\[
                 \operatorname{length}_O\mathcal S_j[b]
                               \le\mathcal M_j.\mathrm{len}.
                                                               \tag{G52.18}
\]
No claim of finiteness has been used in this application.

\proofstep{Pass from bounded torsion to the full group.}
Every continuous cocycle from the compact group $G_K$
to the discrete module $D(\alpha_j)$ has finite image.
A nonzero element of $O$ kills that finite set, so
every cohomology class and every element of
$\mathcal S_j$ is torsion.
Choose a uniformizer $\varpi$ of $O$.  Then
$\mathcal S_j$ is the directed union of
$\mathcal S_j[\varpi^n]$.
The length of a module is the supremum of the lengths
of its finitely generated submodules.  Every finitely
generated submodule of this torsion module lies in
one $\mathcal S_j[\varpi^n]$.
Taking the supremum in (G52.18) proves (G52.9), including
the possible infinite value.
This is exactly the torsion/direct-limit step
\texttt{first\_cohomology\_torsion} and
\texttt{len\_of\_torsion} in
\path{UnitaryThreeGroups/VlenLe.lean}.
\end{proof}

\paragraph{Scope of the reduction.}
For $j=2$, (G52.1) is the desired
$\chi\psi^{-1}\kappa_\ell^{-e}$ coefficient; $j=1$
gives its cubic conjugate.
The length comparison is attached to the constructed
model for that same seed, not to the old untwisted
model or to an assumed equality of deformation rings.
To get a numerical upper bound one must still supply
the actual model's finite-support and point-rigidity
inputs, including the projected mixed tangent
condition and Type I automorphy used in extraction.
Coefficient primes excluded above and split primes
with nontrivial local cubic character remain separate.

\paragraph{Geometric dictionary.}
The induced lattice (G52.10) is the explicit
restriction-of-scalars description of a rank-one
finite-character local system followed by its Tate
twist, in the previously audited SGA~1, Expos\'e~V,
\S\S4 and~7 finite-cover framework.
The new content is the integral normalized deformation
and its cotangent-length comparison.  These are proved
by the displayed cocycle and square-zero arguments and
the expressly retained \texttt{LemmaM'}, not inferred
from the finite-cover dictionary or from an unconstructed
geometric object.
